\section{Definición del negocio y estudio de mercado regional}
\label{sec:mercado-regional}

\subsection{Base de evaluación y alcance del plan}

Este plan amplía el trabajo de la etapa institucional con un estudio de mercado,
la organización del servicio y la proyección financiera. La revisión sigue los
28 indicadores de la rúbrica regional \parencite{rubrica_plan_negocios_2026}.
El apéndice indica dónde se desarrolla cada criterio, incluidos m y o, que
repiten el mismo requisito. Las marcas de tres puntos del formulario no son
una calificación recibida por Pulso.

La estructura sigue los lineamientos de ExpoTÉCNICA \parencite{mep2026}.
Los cálculos propios se organizaron en cinco apartados de trabajo: resumen,
clientes e ingresos, egresos e inversión, equilibrio y estados financieros.
Las tablas de este informe presentan ese escenario preliminar y los supuestos
que lo sostienen.

\subsection{Identidad, misión, visión y actividad}

Pulso es la marca comercial propuesta para un servicio de software por suscripción
orientado a la gestión de gimnasios. Su validación inicial se sitúa en Pérez
Zeledón, Costa Rica. La actividad prevista consiste en desarrollar, mantener y
acompañar el uso de una plataforma web que conecta información administrativa y
seguimiento deportivo. El negocio se encuentra en etapa de proyecto estudiantil
y prototipo; la forma jurídica comercial sigue por definir.

La misión es ayudar a responsables, entrenadores y miembros a organizar las
membresías, los pagos, la asistencia y el entrenamiento desde sus dispositivos
habituales. La visión es mantener un servicio útil para gimnasios locales y,
después, ampliar su alcance regional según la demanda y la capacidad del equipo.

Pulso conecta tareas que algunas personas consultadas realizan con herramientas
separadas. Por ejemplo, el responsable revisa la membresía y el pago, mientras
el miembro consulta su rutina desde su propia pantalla. El acompañamiento y el
acceso móvil completan la propuesta. En la demostración se compararán estas
tareas con el método que ya utiliza cada gimnasio.

\subsection{Comprador, usuario y mercado alcanzable}

El dueño o administrador decide si paga la suscripción. El entrenador organiza
las rutinas y el miembro consulta su información. Por eso interesa conocer
tanto la disposición de compra del responsable como la utilidad que encuentran
quienes usarán la plataforma.

Para la demostración se buscará un gimnasio con internet, un dispositivo
disponible y necesidad de organizar al menos dos procesos. También debe
participar quien decide sobre la contratación y estar dispuesto a revisar los
resultados. La cantidad de miembros ayudará a estimar el soporte y la capacidad
que necesitaría el servicio.

La evidencia disponible consta de 15 respuestas de conveniencia, cinco de
responsables y diez de miembros. Cuatro responsables mencionaron dificultades
con cobros; cuatro solicitaron una demostración antes de indicar un precio.
Nueve miembros usarían principalmente el celular. Estos hallazgos justifican
la demostración administrativa y el acceso móvil, pero no permiten calcular la
demanda total ni una cuota del mercado local
\parencite{encuesta_gymhub_2026}.

Para el escenario preliminar se supone contactar diez posibles clientes y
conseguir cinco contrataciones iniciales: una conversión del 50\%. Es un
supuesto del equipo que aún debe comprobarse. Luego se plantea una contratación
por mes y dos cancelaciones en el año. Así, 16 altas menos dos bajas dejan
14 gimnasios activos al cierre. Se registrarán los contactos, las
contrataciones y las cancelaciones para comparar esta proyección con lo que
ocurra durante la venta.

\subsection{Competencia y decisión de compra}

La comparación de alternativas desarrollada en los antecedentes se aplica a una
decisión concreta: conservar herramientas actuales o asumir una suscripción y un
cambio de procedimiento. Papel y hojas de cálculo ofrecen familiaridad; una
plataforma especializada puede ofrecer funciones ya consolidadas. Pulso debe
mostrar valor en tareas específicas, facilidad de aprendizaje y acompañamiento,
sin afirmar que todas las alternativas sean costosas o inadecuadas.

Antes de contratar, el gimnasio recibirá una demostración con datos ficticios y
una explicación de funciones, precios y límites. El acuerdo de prueba indicará
la duración, el soporte y cómo terminar el servicio. También se probará la
exportación de información antes de ofrecerla. La separación de datos entre
gimnasios y la personalización siguen pendientes para la etapa comercial.

\subsection{Plan de promoción, distribución y permanencia}

La promoción se apoyará en contacto directo, mensajes autorizados, una página
informativa y material audiovisual de tareas reales. La venta se realizará
mediante conversación y demostración con el comprador. La entrega será digital,
a través del acceso web y la configuración guiada. El soporte utilizará correo
o mensajería. Se anotarán las consultas para darles seguimiento y conocer el
tiempo que requiere la atención.

\begin{longtable}{p{2.1cm}p{4.0cm}p{4.0cm}p{3.8cm}}
\caption{Calendario comercial propuesto e indicadores de seguimiento}\\
\toprule
Periodo & Acción y responsable & Recurso presupuestado & Evidencia o medida \\
\midrule\endfirsthead
\toprule Periodo & Acción y responsable & Recurso presupuestado & Evidencia o medida \\
\midrule\endhead
Antes del mes 1 & Equipo: preparar demostración, condiciones y lista de posibles clientes. & Trabajo de desarrollo e investigación ya incluido en el presupuesto. & Ficha de alcance, permisos y registro de contactos. \\
Mes 1 & Equipo: contacto y demostración a diez prospectos como hipótesis de trabajo. & CRC 40\,000 en viáticos, CRC 5\,000 en publicidad y 36 horas comerciales valoradas. & Contactados, demostraciones y cinco altas supuestas; registrar también rechazos. \\
Meses 2--12 & Equipo: seguimiento, incorporación y atención. & CRC 12\,000 mensuales en viáticos, CRC 5\,000 en publicidad y 12 horas comerciales por mes. & Altas, bajas, permanencia, costo por alta y tiempo de soporte. \\
Cada cierre & Equipo: revisar uso y decisión de continuidad. & Soporte valorado en 0,5 horas por gimnasio al mes. & Incidencias resueltas, tareas completadas y motivos de cancelación. \\
\bottomrule
\end{longtable}
{\small\textit{Nota}. Elaboración propia a partir del escenario financiero
preliminar. Las actividades y sus costos están incluidos en la proyección.}

La conversión se medirá como altas entre prospectos contactados; la cancelación,
como bajas entre clientes al inicio del periodo. Para estudiar la permanencia
se seguirá a los clientes que entraron en un mismo mes. Si las demostraciones
no llevan a contrataciones, se revisará la oferta antes de aumentar la
publicidad. El costo de conseguir un cliente incluirá viajes, anuncios y horas
de trabajo.

\subsection{Calidad y pertinencia medibles}

El seguimiento propuesto compara tareas antes y durante una prueba autorizada:
registrar un pago, consultar una vigencia, registrar asistencia y resolver una
rutina. Se anotarán duración, éxito, errores y necesidad de ayuda, conservando
las mismas condiciones de observación. El equipo revisará los incidentes cada
semana y comprobará las correcciones con la persona que los reportó.

La prueba permitirá saber si el gimnasio acepta el precio, si puede completar
las tareas y si entiende las pantallas. También se contará cuántos formularios
impresos se sustituyen. Para medir las mejoras se comparará el proceso habitual
con el uso de Pulso; esta evaluación aún está pendiente.
